\section{Annexe : Schéma des données JSON}
\label{sec:data-schema}

\subsection{data.json --- Dictionnaire}

Le fichier \texttt{static/data.json} contient l'intégralité du dictionnaire, organisé par catégorie grammaticale au premier niveau.

\subsubsection{Structure générale}

\begin{lstlisting}[title=Structure de premier niveau]
{
  "prep":    { ... },    // Prepositions
  "nom":     { ... },    // Noms
  "adj":     { ... },    // Adjectifs
  "card":    { ... },    // Cardinaux
  "proposs": { ... },    // Pronoms possessifs
  "pron":    { ... },    // Pronoms (determinants)
  "proper":  { ... },    // Pronoms personnels
  "verb":    { ... },    // Verbes
  "adv":     { ... },    // Adverbes
  "conj":    { ... },    // Conjonctions
  "part":    { ... }     // Particules
}
\end{lstlisting}

\subsubsection{Format des formes fléchies}

Chaque forme fléchie est stockée sous l'un de ces formats~:

\begin{enumerate}
  \item \textbf{Paire simple} : \texttt{["текст", position]}

  Exemple : \texttt{["будинку", 3]} signifie \ukr{буди́нку} (accent sur la 3\ieme{} lettre).

  \item \textbf{Variantes multiples} : \texttt{[["варіант1", pos1], ["варіант2", pos2]]}

  Pour les mots ayant plusieurs formes acceptées. La sélection se fait via \texttt{var=N} dans \texttt{data-info}.

  \item \textbf{Ancien format (legacy)} : \texttt{["форма1", "форма2", position]}

  Plusieurs formes partageant la même position d'accent. Converti automatiquement par \texttt{toPairs()}.
\end{enumerate}

Conventions de position~:
\begin{itemize}
  \item $> 0$ : position de l'accent (1-indexed).
  \item $-1$ : mot monosyllabique (pas d'accent visible).
  \item $-2$ : accent inconnu.
\end{itemize}

\subsubsection{Noms (\texttt{nom})}

\begin{lstlisting}[language=javascript]
"nom": {
  "LEMME": {
    "cas": {
      "nomi": { "s": ["forme_s", pos], "pl": ["forme_pl", pos] },
      "gen":  { "s": [...], "pl": [...] },
      "dat":  { "s": [...], "pl": [...] },
      "acc":  { "s": [...], "pl": [...] },
      "ins":  { "s": [...], "pl": [...] },
      "loc":  { "s": [...], "pl": [...] }
    },
    "genre": "m" | "f" | "n",
    "base_html": "<span class=\"ukr\" data-info=\"LEMME;nom;cas;nomi;s\">...</span>",
    "phrases": { "phraseKey1": true, ... }   // optionnel
  }
}
\end{lstlisting}

Les 6 cas ukrainiens~:
\begin{center}
\begin{tabular}{lll}
  \toprule
  \textbf{Clé} & \textbf{Cas} & \textbf{Ukrainien} \\
  \midrule
  \texttt{nomi} & Nominatif & \ukr{називний} \\
  \texttt{gen} & Génitif & \ukr{родовий} \\
  \texttt{dat} & Datif & \ukr{давальний} \\
  \texttt{acc} & Accusatif & \ukr{знахідний} \\
  \texttt{ins} & Instrumental & \ukr{орудний} \\
  \texttt{loc} & Locatif & \ukr{місцевий} \\
  \bottomrule
\end{tabular}
\end{center}

\subsubsection{Adjectifs, pronoms possessifs, cardinaux, pronoms (\texttt{adj}, \texttt{proposs}, \texttt{card}, \texttt{pron})}

Même structure que les noms, mais avec 4 genres au lieu de 2 nombres~:

\begin{lstlisting}[language=javascript]
"adj": {
  "LEMME": {
    "cas": {
      "nomi": { "m": [...], "f": [...], "n": [...], "pl": [...] },
      "gen":  { "m": [...], "f": [...], "n": [...], "pl": [...] },
      // ... autres cas
    },
    "base_html": "...",
    "phrases": { ... }
  }
}
\end{lstlisting}

\subsubsection{Pronoms personnels (\texttt{proper})}

Pas de genre dans la déclinaison --- seulement les cas~:

\begin{lstlisting}[language=javascript]
"proper": {
  "LEMME": {
    "cas": {
      "nomi": ["forme", pos],
      "gen":  ["forme", pos],
      // ...
    },
    "base_html": "..."
  }
}
\end{lstlisting}

\subsubsection{Verbes (\texttt{verb})}

\begin{lstlisting}[language=javascript]
"verb": {
  "LEMME": {
    "inf": ["infinitif", pos],
    "conj": {
      "pres": {           // present (verbes imperfectifs)
        "1p": { "s": [...], "pl": [...] },
        "2p": { "s": [...], "pl": [...] },
        "3p": { "s": [...], "pl": [...] }
      },
      "fut": { ... },     // futur (meme structure que pres)
      "pass": {            // passe (structure par genre)
        "m":  { "s": [...], "pl": [...] },
        "f":  { "s": [...], "pl": [...] },
        "n":  { "s": [...], "pl": [...] }
      },
      "imp": {             // imperatif (seulement 1p et 2p)
        "1p": { "s": [...], "pl": [...] },
        "2p": { "s": [...], "pl": [...] }
      },
      "impers": [...]      // forme impersonnelle (optionnel)
    },
    "asp": "perfectif" | "imperfectif",
    "coupl": "LEMME_ASPECT_OPPOSE",  // optionnel
    "base_html": "...",
    "phrases": { ... }
  }
}
\end{lstlisting}

\subsubsection{Mots invariables (\texttt{adv}, \texttt{conj}, \texttt{part}, \texttt{prep})}

\begin{lstlisting}[language=javascript]
"adv": {
  "LEMME": {
    "base": ["forme", pos],
    "base_html": "..."
  }
}
\end{lstlisting}

% ================================================================
\subsection{phrases.json --- Phrases d'exemple}

\begin{lstlisting}[language=javascript]
{
  "CLE_PHRASE": {
    "phrase_html": "<span class=\"ukr\" data-info=\"...\">mot</span> ...",
    "traduction": "Traduction francaise",
    "ref": {
      "Auteur": "L1,Cours5,v1.0"
    },
    "remarque": "Note optionnelle",
    "genereVerbe": {             // optionnel
      "verbe": "LEMME_VERBE",
      "temps": "pres",
      "frag1": "texte avant",
      "frag2": "texte apres"
    }
  }
}
\end{lstlisting}

\begin{description}
  \item[\texttt{phrase\_html}] Le contenu HTML de la phrase, avec des spans \texttt{.ukr} interactifs.
  \item[\texttt{traduction}] Traduction française de la phrase.
  \item[\texttt{ref}] Métadonnées de référence (auteur, cours, source). Utilisées pour le filtrage.
  \item[\texttt{remarque}] Note optionnelle affichée sous la phrase.
  \item[\texttt{genereVerbe}] Si présent, permet de générer toutes les formes conjuguées du verbe pour le temps indiqué, insérées dans le contexte \texttt{frag1 ... pronom ... forme ... frag2}.
\end{description}
